% theme: applications_of_similarity_ratios
\MajorQuestionLesson{相似比の利用}
\SetRowSpace{3mm}

\Needspace{80mm}
\TypeQuestion{}{similarity_application_right_triangle}
\FigProblemBlock{53mm}{fig/lesson_04_right_altitude}{
\Text{図の $\Tri{ABC}$ は $\Ang{BAC}=90^\circ$ の直角三角形である。点 $\mathrm{A}$ から辺 $\Lseg{BC}$ に垂線 $\Lseg{AH}$ を引いた。$\Lseg{BH}=4\mathrm{cm}$、$\Lseg{CH}=9\mathrm{cm}$ である。}
\QQ{$\Tri{ABH}$ と相似な三角形を、対応する頂点の順に2つ書きなさい。}{}
\QQ{$\Lseg{AH}$ の長さを求めなさい。}{}
\QQ{$\Lseg{AB}$ の長さを求めなさい。}{}
}{
\QQ{$\Lseg{AC}$ の長さを求めなさい。}{}
}

\Needspace{80mm}
\TypeQuestion{}{similarity_application_equilateral}
\FigProblemBlock{53mm}{fig/lesson_04_equilateral}{
\Text{図の $\Tri{ABC}$ は1辺が $12\mathrm{cm}$ の正三角形である。$\Lseg{BD}=8\mathrm{cm}$ となる点 $\mathrm{D}$ を辺 $\Lseg{BC}$ 上にとり、$\Lseg{DE}\perp\Lseg{BC}$ となる点 $\mathrm{E}$ を辺 $\Lseg{AC}$ 上にとった。また、$\Lseg{AH}\perp\Lseg{BC}$ とする。}
\QQ{$\Lseg{CH}$ の長さを求めなさい。}{}
\QQ{$\Tri{CDE}$ と相似な三角形を、対応する頂点の順に書きなさい。}{}
}{
\QQ{$\Lseg{CE},\Lseg{AE}$ の長さをそれぞれ求めなさい。}{}
\QQ{$\Lseg{DE}:\Lseg{AH}$ を求めなさい。}{}
}

\Needspace{80mm}
\TypeQuestion{}{similarity_application_isosceles}
\FigProblemBlock{53mm}{fig/lesson_04_isosceles}{
\Text{図の $\Tri{ABC}$ は $\Lseg{AB}=\Lseg{AC}=9\mathrm{cm}$、$\Lseg{BC}=6\mathrm{cm}$ の二等辺三角形である。辺 $\Lseg{AB}$ 上に、$\Ang{BCD}=\Ang{BAC}$ となる点 $\mathrm{D}$ をとった。}
\QQ{$\Tri{BCD}$ と相似な三角形を、対応する頂点の順に書きなさい。}{}
\QQ{$\Lseg{BD}$ の長さを求めなさい。}{}
\QQ{$\Lseg{AD}$ の長さを求めなさい。}{}
}{
\QQ{$\Lseg{CD}$ の長さを求めなさい。}{}
}

\LessonQuestionBreak

\Needspace{80mm}
\TypeQuestion{}{similarity_application_rectangle_fold_diagonal}
\FigProblemBlock{57mm}{fig/lesson_04_fold_diagonal}{
\Text{図の長方形 $\mathrm{ABCD}$ は、$\Lseg{AB}=8\mathrm{cm}$、$\Lseg{AD}=6\mathrm{cm}$、$\Lseg{AC}=10\mathrm{cm}$ である。辺 $\Lseg{AB}$ が対角線 $\Lseg{AC}$ に重なるように折り返すと、点 $\mathrm{B}$ は点 $\mathrm{E}$ に移った。折り目と辺 $\Lseg{BC}$ の交点を $\mathrm{F}$ とする。}
\QQ{$\Lseg{AE},\Lseg{CE}$ の長さをそれぞれ求めなさい。}{}
}{
\QQ{$\Tri{CEF}$ と相似な三角形を、対応する頂点の順に書きなさい。また、相似条件を書きなさい。}{}
\QQ{$\Lseg{CF},\Lseg{BF}$ の長さをそれぞれ求めなさい。}{}
}

\Needspace{80mm}
\TypeQuestion{}{similarity_application_rectangle_fold_opposite}
\FigProblemBlock{57mm}{fig/lesson_04_fold_opposite}{
\Text{図の長方形 $\mathrm{ABCD}$ は、$\Lseg{AB}=13\mathrm{cm}$、$\Lseg{AD}=5\mathrm{cm}$ である。辺 $\Lseg{AB}$ 上の点 $\mathrm{E}$ は $\Lseg{AE}=1\mathrm{cm}$ を満たす。点 $\mathrm{D}$ が点 $\mathrm{E}$ に重なるように折り返すと、折り目は点 $\mathrm{C}$ と辺 $\Lseg{AD}$ 上の点 $\mathrm{F}$ を通った。}
\QQ{$\Lseg{BE},\Lseg{CE}$ の長さをそれぞれ求めなさい。}{}
}{
\QQ{$\Tri{AEF}$ と相似な三角形を、対応する頂点の順に書きなさい。また、相似条件を書きなさい。}{}
\QQ{$\Lseg{AF},\Lseg{DF}$ の長さをそれぞれ求めなさい。}{}
}

\Needspace{80mm}
\TypeQuestion{}{similarity_application_map_scale}
\FigProblemBlock{57mm}{fig/lesson_04_map_scale}{
\Text{縮尺 $1:25000$ の地図上で、地点 $\mathrm{A}$ から地点 $\mathrm{B}$ を通って地点 $\mathrm{C}$ まで進む道を考える。地図上の $\Lseg{AB},\Lseg{BC}$ の長さは図のとおりで、単位は $\mathrm{cm}$ である。ただし、図は略図である。}
\QQ{実際の $\Lseg{AB}$ の道のりを $\mathrm{m}$ で求めなさい。}{}
}{
\QQ{地点 $\mathrm{A}$ から地点 $\mathrm{B}$ を通って地点 $\mathrm{C}$ まで進む実際の道のりを $\mathrm{km}$ で求めなさい。}{}
\QQ{同じ道を縮尺 $1:10000$ の地図に表すと、$\Lseg{AB}$ と $\Lseg{BC}$ を合わせた地図上の長さは何 $\mathrm{cm}$ になるか求めなさい。}{}
}

\Needspace{80mm}
\TypeQuestion{}{similarity_application_river_scale_drawing}
\FigProblemBlock{57mm}{fig/lesson_04_river_scale}{
\Text{川岸の2点 $\mathrm{A},\mathrm{B}$ と、対岸の点 $\mathrm{P}$ の位置を調べ、縮尺 $1:400$ の縮図をかいた。点 $\mathrm{P}$ から直線 $\Lseg{AB}$ に引いた垂線の足を $\mathrm{H}$ とする。図中の長さは縮図上の長さで、単位は $\mathrm{cm}$ である。川の両岸は平行で、川幅は一定とする。}
\QQ{実際の川幅 $\Lseg{PH}$ を $\mathrm{m}$ で求めなさい。}{}
}{
\QQ{実際の $\Lseg{AB}$ の長さを $\mathrm{m}$ で求めなさい。}{}
\QQ{同じ場所を縮尺 $1:1000$ の縮図に表すと、線分 $\Lseg{PH}$ の長さは何 $\mathrm{cm}$ になるか求めなさい。}{}
}
